(1) 质量 \(m\)、正电荷 \(q\) 的非相对论粒子被束缚在半径 \(R\) 的一维圆环上，环内磁场为零，但环围绕磁通 \(\Phi\)。取环向磁矢势 \(A_\theta=\Phi/(2\pi R)\)，写出Hamilton算符。
(2) 求满足单值周期边界条件的能量本征态与能级，定义磁通量子 \(\Phi_0=h/q\)。
(3) 证明能谱集合对 \(\Phi\to\Phi+\Phi_0\) 的周期性，并求半整数磁通处的一个简并能级例子。
(4) 求单个本征态的持续电流 \(I_n=-\partial E_n/\partial\Phi\)，说明为什么局部磁场为零仍可出现可观测的磁通依赖。
